% =============================================================================
%  6. Discussion
%  Every limitation here is one WE measured. Writing them before a reviewer
%  finds them is the point; each paragraph names the number that bounds it.
% =============================================================================
\section{Discussion}
\label{sec:discussion}

\subsection{The boundary of our own claim}

Object recall is pinned at 55/73 and no proposal stage moves it. Eighteen
instances are localized and then rejected, four of them by the feature
\proposer{} alone --- that is, the \proposer{} found them and the \judge{}
said no. This is the exact boundary of the claim this paper makes: the
proposal stage sets the ceiling, and the remaining headroom sits at the
\judge{}, where \tabref{tab:judge} shows prompt and model selection to be
exhausted for this class of local model. The next lever is therefore
adaptation --- parameter-efficient fine-tuning of the \judge{} on human-labelled
crops from this deployment --- which is consistent with the enabling role that
Borodin et al.~\cite{Borodin_2025} report for PEFT in a neighbouring task. We
consider that the natural continuation of this work rather than an afterthought.

\subsection{Threats to validity}

\paragraph{Dataset scale and provenance.}
Sixty-eight cases, 31 of them anomalous, from a single deployment, with eleven
cases carrying inpainted rather than naturally occurring anomalies. Two of the
four anomaly types are represented by six and seven instances respectively.
Nothing in \secref{sec:results:main} should be read as a precise operating
point; what it supports is a comparison \emph{within} the table, where all four
arms see exactly the same 68 cases and the same ground truth. The
public-data evaluation of \secref{sec:results:cdnet} is the mitigation, and it
is a partial one: it validates the component carrying the claim, on 5\,000-plus
annotated frames we did not build, but it cannot validate the pipeline
end to end because CDnet has no notion of a semantic verdict.

\paragraph{The dataset cannot be released.}
The footage is operational CCTV of a public-transport interior and contains
identifiable passengers. This is a real limitation of the work and the reason
every component-level claim in this paper is also measured on public data. It
is not, however, a peculiarity of this project: the application-oriented review
of in-vehicle monitoring by Caetano et al.~\cite{Caetano_2022} reports that
work in this domain remains ``fully dependent on the availability of private
datasets'', and identifies the extension of outdoor abandoned-object detection
to vehicle interiors as an opportunity the accessible literature has
overlooked. Their proposed remedy for the scarcity --- synthetic augmentation
--- is the family our eleven inpainted cases belong to, and we regard building
a releasable in-vehicle benchmark as the field-level work this paper cannot do
on its own.

\paragraph{A single judge family carries the end-to-end numbers.}
\tabref{tab:main} uses one \judge{} across all four arms by construction ---
that is what makes it a controlled comparison --- so the invariance of object
recall is established for that \judge{}. \tabref{tab:judge} shows that the
three alternative model families reach 39--64 instances over a fixed set of
boxes, which bounds but does not eliminate the concern that the invariance is
model-specific.

\paragraph{Statistical power of the per-camera AUROCs.}
Four of the five columns in \tabref{tab:baselines} and
\tabref{tab:dinomaly2} rest on eight image pairs or fewer, where a single
swapped pair moves the AUROC by 0.125. We report them for completeness and
draw conclusions only from the 176-pair column and from the score-level
statistics, which pool many more measurements than the AUROCs do.

\subsection{What the deployment requires}

\paragraph{One model per camera.}
\secref{sec:results:transfer} establishes this as a measured requirement rather
than an implementation choice: a foreign checkpoint is blind rather than
miscalibrated, so no per-camera threshold rescues it, and the authors' own
remedy --- context-aware recentering --- halves the gap while leaving transfer
at chance and costing about 0.10 AUROC on the home camera. For a fleet, the
consequence is a per-camera training step of a few minutes on nominal footage,
which is operationally acceptable but must be planned for. It also means that a
camera whose aim is adjusted needs retraining, not merely rethresholding.

\paragraph{The reference must stay valid.}
This is the single assumption on which everything rests, and CDnet quantifies
what happens when it breaks: \texttt{tramstop} at $F_1$ 0.138 and
\texttt{winterDriveway} at 0.143 are sequences whose reference goes stale under
drifting outdoor light. In a tram the natural refresh point is the end-of-service
inspection, when the vehicle is known to be empty. A system that cannot
guarantee that refresh should expect the proposal stage to degrade toward those
numbers rather than toward the 0.978 of \texttt{streetLight}.

\paragraph{The strongest proposal stage is strongest conditionally.}
The feature \proposer{} wins decisively on our fleet and fails on CDnet
(\secref{sec:results:cdnet}), for a reason we could isolate: it assumes a fixed
camera whose reference shares its grid, and it quantises boxes to a 24~px
patch. That is a good trade when the anomaly spans a substantial fraction of a
patch and a bad one on low-resolution outdoor footage evaluated against
pixel-accurate masks. Practitioners should read our recommendation as
conditional on the deployment geometry, not as a general ranking of proposal
mechanisms.

\subsection{Prompting, stated carefully}

Our sweep shows a prompt moving an under-firing model from 39/73 to 48/73 at
unchanged specificity and moving a calibrated model not at all. We read that as
calibration rather than capability, \emph{at this scale, for these models, and
without adaptation}. We deliberately do not generalise further: VERA
\cite{Ye_2025} obtains real gains by treating prompts as learnable parameters,
which is a different intervention from selecting among hand-written
formulations, and the prompt-sensitivity literature surveyed
in~\cite{Ren_2025} reports the effect varying widely across tasks and model
scales. What our result does support is a practical rule --- measure the YES
rate before attributing a recall difference to a prompt --- and a direction:
once a model is calibrated, further formulation work has no headroom to
exploit.

\subsection{What we could not do}

The event-level protocol of Luna et al.~\cite{Luna_2018} would let us report
per-object precision, recall and $F_1$ on abandoned-object sequences, and is
the natural complement to \secref{sec:results:cdnet}. We were unable to obtain
the temporal annotation package it requires, and ABODA ships none of its own,
which is why \secref{sec:results:aboda} is qualitative. This is a missing
measurement, not a missing capability: the evaluation code exists and the
sequences are public.
